% Blueprint: crt-test
% Created by Numina

\begin{document}

\section{Introduction}

The Chinese Remainder Theorem (CRT) is a classical result in elementary number
theory describing the simultaneous solutions of systems of linear congruences
with pairwise coprime moduli. This blueprint formalizes the basic statements:
the two-moduli case, the general existence theorem for $k$ pairwise coprime
moduli, the uniqueness of the solution modulo the product of the moduli, and
the corresponding unique residue class characterization.

\section{Coprimality}

\begin{definition}[Coprime integers]
    \label{def:coprime}
    Two integers $a, b \in \mathbb{Z}$ are said to be \emph{coprime} if
    $\gcd(a, b) = 1$.
\end{definition}

\begin{definition}[Pairwise coprime family]
    \label{def:pairwise-coprime}
    \uses{def:coprime}
    A finite family of integers $(n_1, n_2, \ldots, n_k)$ with $k \geq 1$ is
    said to be \emph{pairwise coprime} if for every pair of indices
    $i \neq j$ with $1 \leq i, j \leq k$ we have $\gcd(n_i, n_j) = 1$.
\end{definition}

\section{Existence}

\begin{lemma}[Two-moduli CRT]
    \label{lem:crt-two}
    \uses{def:coprime}
    Let $m, n \in \mathbb{Z}$ be coprime, i.e.\ $\gcd(m, n) = 1$. For any
    $a, b \in \mathbb{Z}$ there exists an integer $x \in \mathbb{Z}$ such that
    \[
        x \equiv a \pmod{m}, \qquad x \equiv b \pmod{n}.
    \]
\end{lemma}
\begin{proof}
    \uses{def:coprime}
    Since $\gcd(m, n) = 1$, by Bézout's identity there exist integers
    $u, v \in \mathbb{Z}$ with $u m + v n = 1$. Set
    \[
        x \;=\; b \cdot u m \;+\; a \cdot v n.
    \]
    Reducing modulo $m$ gives $x \equiv a \cdot v n \equiv a \cdot 1 \equiv a
    \pmod m$, and reducing modulo $n$ gives $x \equiv b \cdot u m \equiv b
    \cdot 1 \equiv b \pmod n$.
\end{proof}

\begin{theorem}[General CRT existence]
    \label{thm:crt-existence}
    \uses{def:pairwise-coprime, lem:crt-two}
    Let $k \geq 1$ and let $n_1, n_2, \ldots, n_k$ be pairwise coprime
    integers. For any integers $a_1, a_2, \ldots, a_k$ there exists
    $x \in \mathbb{Z}$ such that
    \[
        x \equiv a_i \pmod{n_i} \quad \text{for every } 1 \leq i \leq k.
    \]
\end{theorem}
\begin{proof}
    \uses{lem:crt-two, def:pairwise-coprime}
    Induction on $k$. The case $k = 1$ is trivial: take $x = a_1$. Suppose
    the result holds for $k$ pairwise coprime moduli, and consider $k+1$
    pairwise coprime moduli $n_1, \ldots, n_{k+1}$. By the induction
    hypothesis there exists $y \in \mathbb{Z}$ with $y \equiv a_i \pmod{n_i}$
    for $1 \leq i \leq k$. Let $M = n_1 n_2 \cdots n_k$. Since each $n_i$ is
    coprime to $n_{k+1}$, the product $M$ is also coprime to $n_{k+1}$. By
    the two-moduli CRT (Lemma~\ref{lem:crt-two}) applied to the moduli $M$
    and $n_{k+1}$, there exists $x \in \mathbb{Z}$ with $x \equiv y \pmod M$
    and $x \equiv a_{k+1} \pmod{n_{k+1}}$. Then for each $1 \leq i \leq k$,
    $n_i \mid M$ and $x \equiv y \equiv a_i \pmod{n_i}$, while
    $x \equiv a_{k+1} \pmod{n_{k+1}}$ by construction.
\end{proof}

\section{Uniqueness}

\begin{theorem}[CRT uniqueness]
    \label{thm:crt-uniqueness}
    \uses{def:pairwise-coprime}
    Let $n_1, \ldots, n_k$ be pairwise coprime integers and set
    $N = n_1 n_2 \cdots n_k$. If $x, y \in \mathbb{Z}$ satisfy
    $x \equiv y \pmod{n_i}$ for every $1 \leq i \leq k$, then
    $x \equiv y \pmod{N}$.
\end{theorem}
\begin{proof}
    \uses{def:pairwise-coprime}
    By hypothesis $n_i \mid (x - y)$ for every $i$. Since the $n_i$ are
    pairwise coprime, their product $N = \prod_{i=1}^{k} n_i$ also divides
    $x - y$ (a standard consequence of the fact that if $a$ and $b$ are
    coprime and both divide $c$, then $ab \mid c$, applied iteratively).
    Hence $x \equiv y \pmod{N}$.
\end{proof}

\begin{corollary}[Unique residue class form of CRT]
    \label{cor:crt-unique-class}
    \uses{thm:crt-existence, thm:crt-uniqueness}
    Let $n_1, \ldots, n_k$ be pairwise coprime integers, set
    $N = n_1 n_2 \cdots n_k$, and let $a_1, \ldots, a_k \in \mathbb{Z}$.
    Then there exists a unique residue class $\overline{x} \in \mathbb{Z}/N\mathbb{Z}$
    such that every representative $x \in \overline{x}$ satisfies
    $x \equiv a_i \pmod{n_i}$ for every $1 \leq i \leq k$.
\end{corollary}
\begin{proof}
    \uses{thm:crt-existence, thm:crt-uniqueness}
    Existence of some $x \in \mathbb{Z}$ with $x \equiv a_i \pmod{n_i}$ for
    every $i$ is given by Theorem~\ref{thm:crt-existence}. If $x'$ is another
    such integer, then $x \equiv x' \pmod{n_i}$ for every $i$, hence by
    Theorem~\ref{thm:crt-uniqueness} we have $x \equiv x' \pmod N$. Thus the
    residue class $\overline{x} \in \mathbb{Z}/N\mathbb{Z}$ is uniquely
    determined by the data $(a_1, \ldots, a_k)$.
\end{proof}

\end{document}
